\chapter{Thermal Gauge Theory And Screening}
\label{ch:thermal-gauge-theory-screening}

\ConventionsMixed

This chapter studies gauge theories at nonzero temperature through static
correlators and their long-distance screening behavior.  The coupling and
trace conventions are those of the gauge chapters: the Yang--Mills action is
written with
\[
  \frac{1}{4g^2}\operatorname{tr}(F_{\mu\nu}F^{\mu\nu}),
\]
and group invariants are computed using the same invariant bilinear form.

\section{Static Correlators}

Let the Euclidean time circle have length \(\beta=1/T\).  A bosonic field has
Matsubara frequencies \(\omega_n=2\pi nT\).  Static screening is controlled by
the \(n=0\) sector of gauge-invariant correlators.  For an operator
\(\mathcal O(\tau,\mathbf x)\), define its static projection by
\[
  \mathcal O_0(\mathbf x)=T\int_0^\beta d\tau\,\mathcal O(\tau,\mathbf x).
\]
The screening masses are inverse correlation lengths extracted from
\[
  \left\langle \mathcal O_0(\mathbf x)\mathcal O_0(0)\right\rangle_c
  \sim
  \sum_\alpha C_\alpha
  \frac{e^{-M_\alpha |\mathbf x|}}{|\mathbf x|^{(d-2)/2}}
\]
when the static transfer matrix has isolated lowest eigenvalues in the chosen
channel.

\section{Electric Screening}

In a weakly coupled high-temperature regime, the temporal gauge field \(A_0\)
behaves in the static sector as an adjoint scalar.  The electric screening
mass is defined perturbatively by the pole of the static longitudinal
propagator, or nonperturbatively by the lowest gauge-invariant state in the
electric channel after the relevant operator has been specified.

Choose a basis \(T^a_R\) for each representation \(R\) such that
\[
  \operatorname{tr}_R(T^a_RT^b_R)=T_R\delta^{ab},
  \qquad
  f^{acd}f^{bcd}=C_A\delta^{ab}.
\]
With the same coupling \(g\) in the covariant derivative and in the kinetic
term convention above, the one-loop static polarization gives
\[
  m_D^2
  =
  g^2T^2\left(
  \frac{C_A}{3}
  +\frac{1}{3}\sum_{\text{Dirac }f}T_{R_f}
  +\frac{1}{6}\sum_{\text{complex }s}T_{R_s}
  \right).
\]
For \(SU(N)\) with the fundamental trace convention
\(\operatorname{tr}_{\mathbf N}(T^aT^b)=\delta^{ab}\), this means
\(C_A=2N\) and \(T_{\mathbf N}=1\).

\begin{proposition}[Origin of the Debye mass term]
\label{prop:debye-mass-origin}
At one loop, the static quadratic term for \(A_0\) in the dimensionally
reduced action is the zero-frequency, zero-momentum limit of the Euclidean
current-current correlator coupled to \(A_0\).
\end{proposition}

\begin{proof}
Introduce a background temporal field \(A_0^a\) that is independent of
Euclidean time and space.  Expanding the finite-temperature generating
functional to second order gives
\[
  \Delta W^{(2)}
  =
  \frac12 A_0^a A_0^b
  \int_0^\beta d\tau\int d^{d-1}x\,
  \left\langle J_0^a(\tau,\mathbf x)J_0^b(0,0)\right\rangle_c .
\]
Translation invariance identifies the integral with the static polarization
\(\Pi_{00}^{ab}(0,\mathbf0)\).  Gauge symmetry in the unbroken thermal state
sets this tensor proportional to \(\delta^{ab}\).  The coefficient is the
mass matrix of the adjoint scalar \(A_0\) in the static effective action,
which is \(m_D^2\delta^{ab}\).
\end{proof}

\section{Magnetic Sector}

The spatial gauge fields \(A_i\) have no gauge-invariant local mass term in
the dimensionally reduced action.  The high-temperature static magnetic
sector is therefore a three-dimensional Yang--Mills theory with coupling
\[
  g_3^2=g^2T
\]
in four spacetime dimensions.  Its infrared scale is of order \(g_3^2\), and
the corresponding magnetic screening lengths are nonperturbative objects of
the three-dimensional gauge theory.  Perturbative expansions of static
magnetic quantities encounter this scale because massless spatial gauge
fields remain in the reduced action.

\section{Polyakov Loop And Center Symmetry}

The thermal Wilson line at spatial point \(\mathbf x\) is
\[
  P(\mathbf x)=
  \mathcal P\exp\!\left(\ii\int_0^\beta d\tau\,A_0(\tau,\mathbf x)\right).
\]
In a pure gauge theory, its trace in a representation charged under the center
transforms under the center one-form symmetry wrapped on the thermal circle.
The free energy of an external static charge is defined by
\[
  e^{-\beta F_Q}=\left\langle \operatorname{tr}_R P(\mathbf x)\right\rangle
\]
after renormalizing the line.  Correlators of two Polyakov loops define
static color-singlet free energies; their large-distance decay belongs to the
screening spectrum rather than to a zero-temperature particle spectrum.

\section{Dimensional Reduction}

At distances much larger than \(\beta\), nonzero Matsubara modes can be
integrated out into a three-dimensional effective theory.  In four spacetime
dimensions the bosonic part begins as
\[
  S_{\mathrm{EQCD}}
  =
  \int d^3x\,
  \left[
  \frac{1}{4g_3^2}\operatorname{tr}(F_{ij}F_{ij})
  +\frac{1}{2g_3^2}\operatorname{tr}(D_iA_0D_iA_0)
  +\frac{m_E^2}{2g_3^2}\operatorname{tr}(A_0^2)
  +\lambda_E\,\operatorname{tr}(A_0^4)+\cdots
  \right].
\]
The ellipsis denotes higher local operators ordered by the static derivative
and field expansion.  A quantitative use of this action requires matching
the coefficients \(g_3^2,m_E^2,\lambda_E,\ldots\) in the same regulator and
operator scheme as the original thermal theory.

\begin{openproblem}[Gauge-invariant screening spectrum]
\label{op:gauge-invariant-screening-spectrum}
Give a continuum construction of the full static screening spectrum of
thermal non-Abelian gauge theory from gauge-invariant operators, including
the relation between perturbative electric matching, nonperturbative magnetic
dynamics, and Polyakov-loop sectors.
\end{openproblem}
